% !TEX program = xelatex
% ============================================================================
%  POSN Computer Qualification — Polynomials
%  Topic: พหุนาม (Polynomials)
% ============================================================================

\documentclass[11pt, aspectratio=169]{beamer}
% ============================================================================
%  POSN Computer Qualification — Shared Beamer Preamble
%  Used by all topic slide decks. Compile with XeLaTeX.
% ============================================================================

% ---------- Thai language & font ----------
\usepackage{iftex}
\usepackage{fontspec}

% XeTeX-specific: enable Thai line breaking
\ifXeTeX
  \XeTeXlinebreaklocale{th}
\fi
% Noto Sans Thai — body text (clean modern sans-serif, handles all Thai)
\setmainfont{Noto Sans Thai}[
  UprightFont = Noto Sans Thai Light,
  BoldFont    = Noto Sans Thai Medium,
  Scale       = 0.9
]
\setsansfont{Noto Sans Thai}[
  UprightFont = Noto Sans Thai Light,
  BoldFont    = Noto Sans Thai Medium,
  Scale       = 0.9
]

% Noto Sans Thai — clean modern sans-serif for structural elements
\newfontfamily\headingfont{Noto Sans Thai}[
  UprightFont = Noto Sans Thai Medium,
  BoldFont    = Noto Sans Thai Bold,
  Scale       = 0.9
]

% --- Heading font for structural elements ---
\setbeamerfont{title}{family=\headingfont, series=\bfseries, size=\huge}
\setbeamerfont{subtitle}{family=\headingfont, series=\mdseries}
\setbeamerfont{author}{family=\headingfont, series=\mdseries}

% --- Custom Title Page ---
\defbeamertemplate{title page}{custom}[1][]{%
  \centering
  \vspace*{2cm}
  % Title — in white
  {\usebeamerfont{title}\color{white}\inserttitle\par}
  \vspace{0.6cm}
  % Subtitle — in white
  {\usebeamerfont{subtitle}\color{white}\insertsubtitle\par}
  \vspace{2.5cm}
  % Author — blue filled rounded box, white text
  \begin{center}
    \begin{tcolorbox}[arc=3mm, colback=boxBlue, colframe=boxBlue!80,
                      width=0.42\textwidth, halign=center,
                      left=1.5em, right=1.5em, top=0.8em, bottom=0.8em,
                      boxrule=1.5pt, coltext=white]
      \usebeamerfont{author}\insertauthor
    \end{tcolorbox}
  \end{center}
}
\setbeamertemplate{title page}[custom]
\setbeamerfont{frametitle}{family=\headingfont, series=\bfseries}
\setbeamerfont{section in toc}{family=\headingfont}
\setbeamerfont{page number in head/foot}{family=\headingfont}
\setbeamerfont{section title}{family=\headingfont, series=\bfseries}
\setbeamerfont{subsection title}{family=\headingfont}

% Monospace font (for code/verbatim)
\setmonofont{Latin Modern Mono}[Scale=0.9]

% ---------- Math ----------
\usepackage{amsmath, amssymb}

% ---------- Graphics ----------
\usepackage{graphicx}
\usepackage{tikz}

% ---------- String parsing (for section headers) ----------
\usepackage{xstring}

% ---------- Colored boxes ----------
\usepackage[skins, breakable, xparse]{tcolorbox}
% Note: we do NOT load the 'most' option — it predefines example/solution environments
% that would conflict with our custom boxes.

% --- Color definitions ---
\definecolor{boxBlue}{RGB}{41, 128, 185}
\definecolor{boxGreen}{RGB}{39, 174, 96}
\definecolor{boxOrange}{RGB}{230, 126, 34}
\definecolor{boxPurple}{RGB}{142, 68, 173}
\definecolor{boxBlueBg}{RGB}{235, 245, 255}
\definecolor{boxGreenBg}{RGB}{235, 255, 240}
\definecolor{boxOrangeBg}{RGB}{255, 245, 235}
\definecolor{boxPurpleBg}{RGB}{248, 240, 255}

% --- Explanation box (blue) — for definitions, concepts, notes ---
\newtcolorbox{explanation}[1][]{
  enhanced,
  breakable,
  arc         = 2mm,
  colback     = boxBlueBg,
  colframe    = boxBlue!50,
  title       = #1,
  fonttitle   = \bfseries\color{white},
  coltitle    = white,
  colbacktitle= boxBlue,
  attach boxed title to top left = {yshift=-2mm, xshift=2mm},
  boxed title style = {arc=2mm, size=small},
  before skip = 8pt,
  after skip  = 8pt
}

% --- Example box (green) — for worked examples ---
\newtcolorbox{exambox}[1][]{
  enhanced,
  breakable,
  arc         = 2mm,
  colback     = boxGreenBg,
  colframe    = boxGreen!50,
  title       = #1,
  fonttitle   = \bfseries\color{white},
  coltitle    = white,
  colbacktitle= boxGreen,
  attach boxed title to top left = {yshift=-2mm, xshift=2mm},
  boxed title style = {arc=2mm, size=small},
  before skip = 8pt,
  after skip  = 8pt
}

% --- Solution box (orange) — for solutions and answers ---
\newtcolorbox{solbox}[1][]{
  enhanced,
  breakable,
  arc         = 2mm,
  colback     = boxOrangeBg,
  colframe    = boxOrange!50,
  title       = #1,
  fonttitle   = \bfseries\color{white},
  coltitle    = white,
  colbacktitle= boxOrange,
  attach boxed title to top left = {yshift=-2mm, xshift=2mm},
  boxed title style = {arc=2mm, size=small},
  before skip = 8pt,
  after skip  = 8pt
}

% --- Intuition box (purple) — for plain-language explanations, analogies ---
\newtcolorbox{intuition}[1][]{
  enhanced,
  breakable,
  arc         = 2mm,
  colback     = boxPurpleBg,
  colframe    = boxPurple!50,
  title       = #1,
  fonttitle   = \bfseries\color{white},
  coltitle    = white,
  colbacktitle= boxPurple,
  attach boxed title to top left = {yshift=-2mm, xshift=2mm},
  boxed title style = {arc=2mm, size=small},
  before skip = 8pt,
  after skip  = 8pt
}

% ---------- Beamer theme ----------
\usetheme{Madrid}
\usecolortheme{default}

% --- Custom beamer colors (blue accent matching explanation boxes) ---
\setbeamercolor{structure}{fg=boxBlue}
\setbeamercolor{palette primary}{fg=white, bg=boxBlue}
\setbeamercolor{palette secondary}{fg=white, bg=boxBlue!80}
\setbeamercolor{palette tertiary}{fg=white, bg=boxBlue!60}
\setbeamercolor{block title}{fg=white, bg=boxBlue}
\setbeamercolor{block body}{fg=black, bg=boxBlueBg}

% Remove navigation symbols for clean look
\setbeamertemplate{navigation symbols}{}

% Note: Frame title prefix removed — \addtobeamertemplate conflicts with beamer internals

% --- Outline (Table of Contents) styling ---
\setbeamerfont{section in toc}{family=\headingfont, series=\mdseries}
\setbeamerfont{subsection in toc}{family=\headingfont, series=\mdseries}
\setbeamertemplate{section in toc}{%
  \leavevmode\leftskip=1.5em
  \textcolor{boxBlue}{\textbf{\large\inserttocsectionnumber.}}%
  \quad\inserttocsection\par
  \vspace{0.2em}
}
\setbeamertemplate{subsection in toc}{%
  \leavevmode\leftskip=4em
  \textcolor{boxBlue!60}{\small$\bullet$}%
  \quad\inserttocsubsection\par
  \vspace{0.1em}
}

% Section header slides — Thai in white, English in smaller light blue below
\AtBeginSection{
  \begin{frame}[plain]{}
    \vspace*{2cm}
    \begin{beamercolorbox}[sep=12pt, center, rounded=true, shadow=true]{palette primary}
      \IfSubStr{\secname}{(}{%
        \StrBefore{\secname}{(}[\sectionthai]%
        \StrBetween[1,1]{\secname}{(}{)}[\sectioneng]%
        \begin{tabular}{@{}c@{}}
          {\usebeamerfont{title}\sectionthai} \\[0.2em]
          {\usebeamerfont{subtitle}\color{boxBlue!60}\sectioneng}
        \end{tabular}%
      }{%
        {\usebeamerfont{title}\secname\par}%
      }
    \end{beamercolorbox}
    \vfill
  \end{frame}
}

% Subsection header slides — smaller, centered
\AtBeginSubsection{
  \begin{frame}[plain]{}
    \vfill
    \centering
    {\usebeamerfont{section title}\insertsubsectionhead\par}
    \vfill
  \end{frame}
}

% Minimal footer: page number only, right-aligned
\setbeamertemplate{footline}{
  \hfill\usebeamerfont{page number in head/foot}\insertframenumber{} / \inserttotalframenumber\hspace*{1em}\vspace*{0.5ex}
}

% ---------- Spacing & readability ----------
\linespread{1.2}

% ---------- Useful commands ----------

% Highlight important text in blue
\newcommand{\highlight}[1]{\textcolor{boxBlue}{\textbf{#1}}}

% Math set notation helpers
\newcommand{\R}{\mathbb{R}}
\newcommand{\N}{\mathbb{N}}
\newcommand{\Z}{\mathbb{Z}}
\newcommand{\Q}{\mathbb{Q}}

% ============================================================================
%  End of preamble
% ============================================================================


\title[พหุนาม]{พหุนาม}
\subtitle{Polynomials — การดำเนินการ, การหาร, ทฤษฎีบทเศษเหลือ, การแยกตัวประกอบ, ทฤษฎีบททวินาม}
\author[None W. (Yuan)]{None W. (Yuan)}
\date{}

\begin{document}

\begin{frame}[plain]\titlepage\end{frame}
\begin{frame}{สารบัญ (Outline)}\tableofcontents\end{frame}

\begin{frame}{ก่อนเรียน (Prerequisites)}
  \begin{explanation}[สิ่งที่ควรรู้ก่อนเรียน]
    \begin{itemize}
      \item \textbf{จำนวนจริง} — การบวก ลบ คูณ หาร, สมบัติการแจกแจง
      \item \textbf{เลขยกกำลัง} — $x^a \cdot x^b = x^{a+b}$, $(x^a)^b = x^{ab}$
      \item \textbf{สมการ} — การแก้สมการเชิงเส้นและสมการกำลังสอง
    \end{itemize}
  \end{explanation}

  \begin{intuition}[พหุนามคืออะไร?]
    \textbf{พหุนาม} คือนิพจน์ที่ประกอบด้วยตัวแปรยกกำลังด้วยจำนวนเต็มที่ไม่เป็นลบ
    คูณกับค่าคงที่ แล้วบวกกัน เช่น $3x^3 - 5x^2 + 2x - 7$
  \end{intuition}
\end{frame}

% ===========================================================================
%  SECTION 1: นิยามและพื้นฐาน
% ===========================================================================
\section{นิยามและพื้นฐาน (Definition \& Basics)}

\begin{frame}{พหุนามคืออะไร?}
  \begin{explanation}[นิยาม: พหุนามตัวแปรเดียว (Polynomial in One Variable)]
    \textbf{พหุนาม} คือนิพจน์ที่อยู่ในรูป
    \[
      P(x) = a_n x^n + a_{n-1} x^{n-1} + \cdots + a_1 x + a_0
    \]
    โดย $a_n, a_{n-1}, \ldots, a_0 \in \R$ (สัมประสิทธิ์), $a_n \neq 0$, และ $n$ เป็นจำนวนเต็มที่ไม่เป็นลบ
  \end{explanation}

  \begin{columns}[T]
    \column{0.5\textwidth}
      \begin{explanation}[คำศัพท์สำคัญ]
        \begin{itemize}
          \item \textbf{ดีกรี (Degree)}: เลขชี้กำลังสูงสุด = $n$
          \item \textbf{สัมประสิทธิ์นำ}: $a_n$
          \item \textbf{พจน์คงที่}: $a_0$
        \end{itemize}
      \end{explanation}
    \column{0.5\textwidth}
      \begin{exambox}[ตัวอย่าง]
        $P(x) = 4x^5 - 3x^2 + 7x - 1$

        ดีกรี $= 5$, สัมประสิทธิ์นำ $= 4$, พจน์คงที่ $= -1$
      \end{exambox}
  \end{columns}
\end{frame}

\begin{frame}{สิ่งที่ไม่ใช่พหุนาม}
  \begin{explanation}[ข้อควรระวัง]
    นิพจน์ต่อไปนี้ \textbf{ไม่ใช่} พหุนาม:
    \begin{itemize}
      \item $x^{-2} + 3x$ \quad — เลขชี้กำลังติดลบ
      \item $2x^{1/2} + 5$ \quad — เลขชี้กำลังไม่เป็นจำนวนเต็ม
      \item $\dfrac{1}{x} + x^2$ \quad — $x^{-1}$ คือเลขชี้กำลังติดลบ
      \item $3^x + x$ \quad — ตัวแปรเป็นเลขชี้กำลัง (exponential)
    \end{itemize}
  \end{explanation}

  \begin{intuition}[จำง่าย ๆ]
    เลขชี้กำลังของ $x$ ต้องเป็น \highlight{จำนวนเต็มที่ไม่เป็นลบ} ($0, 1, 2, 3, \ldots$) เท่านั้น!
  \end{intuition}
\end{frame}

\begin{frame}{การจำแนกพหุนามตามจำนวนพจน์}
  \begin{columns}[T]
    \column{0.33\textwidth}
      \centering
      \textbf{เอกนาม (Monomial)}

      $5x^3$, \quad $-2x$, \quad $7$

      1 พจน์

    \column{0.33\textwidth}
      \centering
      \textbf{ทวินาม (Binomial)}

      $x^2 + 4$, \quad $3x^3 - 2x$

      2 พจน์

    \column{0.33\textwidth}
      \centering
      \textbf{ไตรนาม (Trinomial)}

      $x^2 + 5x + 6$

      $2x^3 - x^2 + 3$

      3 พจน์
  \end{columns}

  \medskip
  \begin{explanation}[การจำแนกตามดีกรี]
    \begin{itemize}
      \item ดีกรี 0: \textbf{ค่าคงที่} (constant) — เช่น $5$, $-3$
      \item ดีกรี 1: \textbf{เชิงเส้น} (linear) — เช่น $2x + 1$
      \item ดีกรี 2: \textbf{กำลังสอง} (quadratic) — เช่น $x^2 - 4x + 3$
      \item ดีกรี 3: \textbf{กำลังสาม} (cubic) — เช่น $x^3 + 2x^2 - x + 5$
    \end{itemize}
  \end{explanation}
\end{frame}

% ===========================================================================
%  SECTION 2: การดำเนินการของพหุนาม
% ===========================================================================
\section{การดำเนินการของพหุนาม (Polynomial Operations)}

\begin{frame}{การบวกและการลบพหุนาม}
  \begin{explanation}[วิธีทำ]
    \textbf{บวก/ลบพจน์ที่เหมือนกัน} (พจน์ที่มีตัวแปรและเลขชี้กำลังเดียวกัน)
    — บวกหรือลบสัมประสิทธิ์ของพจน์ที่เหมือนกัน
  \end{explanation}

  \begin{exambox}[ตัวอย่าง]
    \textbf{การบวก:} $(3x^3 - 5x^2 + 2x - 1) + (2x^3 + 4x^2 - 3x + 6)
    = 5x^3 - x^2 - x + 5$

    \medskip
    \textbf{การลบ:} $(4x^3 + 2x - 5) - (x^3 - 3x^2 + x + 2)
    = 3x^3 + 3x^2 + x - 7$
  \end{exambox}
\end{frame}

\begin{frame}{การคูณพหุนาม}
  \begin{explanation}[วิธีทำ]
    \textbf{ใช้สมบัติการแจกแจง (Distributive Property)}:
    ทุกพจน์คูณกับทุกพจน์ของอีกพหุนาม แล้วรวมพจน์ที่เหมือนกัน
  \end{explanation}

  \begin{exambox}[ตัวอย่าง]
    \textbf{เอกนาม $\times$ พหุนาม:}
    $3x^2(2x^3 - 5x + 4) = 6x^5 - 15x^3 + 12x^2$

    \medskip
    \textbf{ทวินาม $\times$ ทวินาม:}
    $(2x+3)(x^2-4x+5) = 2x^3 -8x^2 +10x +3x^2 -12x +15
    = 2x^3 -5x^2 -2x +15$
  \end{exambox}
\end{frame}

% ===========================================================================
%  SECTION 3: การหารพหุนาม
% ===========================================================================
\section{การหารพหุนาม (Polynomial Division)}

\begin{frame}{การหารยาว (Long Division)}
  \begin{explanation}[แนวคิด]
    หารพจน์แรก $\rightarrow$ คูณกลับ $\rightarrow$ ลบ $\rightarrow$ ดึงลงมา $\rightarrow$ ทำซ้ำ
  \end{explanation}

  \begin{columns}[T]
    \column{0.55\textwidth}
        \begin{exambox}[ตัวอย่าง: การหารยาว]
        \[
        \begin{array}{rl}
          & x^2 - 2x - 2 \\
          \cline{2-2}
          x - 2 \;\big)\; & x^3 - 4x^2 + 2x + 5 \\
          & -(x^3 - 2x^2) \\
          \cline{2-2}
          & -2x^2 + 2x \\
          & -(-2x^2 + 4x) \\
          \cline{2-2}
          & -2x + 5 \\
          & -(-2x + 4) \\
          \cline{2-2}
          & 1
        \end{array}
        \]
      \end{exambox}

    \column{0.45\textwidth}
      \highlight{ผลหาร $= x^2 - 2x - 2$ \\ เศษ $= 1$}

      \begin{intuition}[ตรวจคำตอบ]
        $(x-2)(x^2-2x-2) + 1$

        $= x^3 - 2x^2 - 2x - 2x^2 + 4x + 4 + 1$

        $= x^3 - 4x^2 + 2x + 5$ \quad $\checkmark$
      \end{intuition}
  \end{columns}
\end{frame}

\begin{frame}{การหารสังเคราะห์ (Synthetic Division)}
  \begin{explanation}[ใช้เมื่อตัวหารอยู่ในรูป $x - c$]
    วิธีลัดที่ใช้เฉพาะสัมประสิทธิ์ — เร็วกว่าการหารยาวเมื่อตัวหารเป็นเชิงเส้น
  \end{explanation}

  \begin{exambox}[ตัวอย่างเดิม: การหารสังเคราะห์]
    \[
    \begin{array}{c|cccc}
      2 & 1 & -4 & 2 & 5 \\
        &   & 2  & -4 & -4 \\
      \hline
        & 1 & -2 & -2 & \boxed{1}
    \end{array}
    \]

    \textbf{ขั้นตอน:} ดึงลงมา $\rightarrow$ คูณด้วย $c$ $\rightarrow$ บวก $\rightarrow$ ทำซ้ำ

    \begin{center}
    {\small $1 \xrightarrow{\times 2} 2 \xrightarrow{+(-4)} -2 \xrightarrow{\times 2} -4 \xrightarrow{+2} -2 \xrightarrow{\times 2} -4 \xrightarrow{+5} 1$}
    \end{center}

    \highlight{ผลหาร $= x^2 - 2x - 2$ เศษ $= 1$}
  \end{exambox}
\end{frame}

% ===========================================================================
%  SECTION 4: ทฤษฎีบทเศษเหลือ
% ===========================================================================
\section{ทฤษฎีบทเศษเหลือ (Remainder Theorem)}

\begin{frame}{ทฤษฎีบทเศษเหลือ}
  \begin{explanation}[ทฤษฎีบท]
    เมื่อหารพหุนาม $P(x)$ ด้วย $(x - c)$ \highlight{เศษที่ได้เท่ากับ $P(c)$}
  \end{explanation}

  \begin{exambox}[ตัวอย่าง]
    $P(x) = x^3 - 4x^2 + 2x + 5$ หารด้วย $(x - 2)$

    \textbf{วิธีที่ 1:} หารยาว/หารสังเคราะห์ $\rightarrow$ เศษ $= 1$

    \textbf{วิธีที่ 2 (ทฤษฎีบทเศษเหลือ):}
    $P(2) = 2^3 - 4(4) + 2(2) + 5 = 8 - 16 + 4 + 5 = 1$

    \highlight{ได้เศษ $= 1$ เท่ากัน!}
  \end{exambox}

  \begin{intuition}[ประโยชน์]
    ไม่ต้องหารจริง! แค่แทนค่า $c$ ลงใน $P(x)$ ก็รู้เศษทันที
  \end{intuition}
\end{frame}

\begin{frame}{ตัวอย่าง: ทฤษฎีบทเศษเหลือ}
  \begin{exambox}[โจทย์]
    จงหาเศษจากการหาร $P(x) = x^4 - 3x^3 + 2x^2 - x + 7$ ด้วย $x + 1$
  \end{exambox}

  \begin{solbox}[วิธีทำ]
    $x + 1 = x - (-1)$ ดังนั้น $c = -1$

    $P(-1) = (-1)^4 - 3(-1)^3 + 2(-1)^2 - (-1) + 7$

    $= 1 - 3(-1) + 2(1) + 1 + 7$

    $= 1 + 3 + 2 + 1 + 7 = 14$

    \highlight{เศษ $= 14$}
  \end{solbox}
\end{frame}

% ===========================================================================
%  SECTION 5: ทฤษฎีบทตัวประกอบ
% ===========================================================================
\section{ทฤษฎีบทตัวประกอบ (Factor Theorem)}

\begin{frame}{ทฤษฎีบทตัวประกอบ}
  \begin{explanation}[ทฤษฎีบท]
    $(x - c)$ เป็น \highlight{ตัวประกอบ} ของ $P(x)$ \textbf{ก็ต่อเมื่อ} $P(c) = 0$
  \end{explanation}

  \begin{columns}[T]
    \column{0.5\textwidth}
      \begin{exambox}[ตัวอย่าง: เป็นตัวประกอบ]
        $P(x) = x^3 - 6x^2 + 11x - 6$

        $P(1) = 1 - 6 + 11 - 6 = 0$

        ดังนั้น $(x-1)$ เป็นตัวประกอบ $\checkmark$
      \end{exambox}
    \column{0.5\textwidth}
      \begin{exambox}[ตัวอย่าง: ไม่เป็นตัวประกอบ]
        $P(x) = x^3 - 2x^2 + x - 3$

        $P(2) = 8 - 8 + 2 - 3 = -1 \neq 0$

        ดังนั้น $(x-2)$ ไม่เป็นตัวประกอบ
      \end{exambox}
  \end{columns}

  \begin{intuition}[ความสัมพันธ์กับทฤษฎีบทเศษเหลือ]
    ถ้าเศษ $= 0$ แปลว่าหารลงตัว $\rightarrow$ ตัวหารเป็นตัวประกอบ
  \end{intuition}
\end{frame}

\begin{frame}{การใช้ทฤษฎีบทตัวประกอบหาค่า $k$}
  \begin{exambox}[โจทย์]
    กำหนด $P(x) = x^3 + kx^2 - 4x + 3$ ถ้า $(x+1)$ เป็นตัวประกอบของ $P(x)$
    จงหาค่าของ $k$
  \end{exambox}

  \begin{solbox}[วิธีทำ]
    $(x+1)$ เป็นตัวประกอบ ดังนั้น $P(-1) = 0$

    $P(-1) = (-1)^3 + k(-1)^2 - 4(-1) + 3$

    $= -1 + k + 4 + 3 = 0$

    $k + 6 = 0$

    \highlight{$k = -6$}
  \end{solbox}
\end{frame}

% ===========================================================================
%  SECTION 6: การแยกตัวประกอบ
% ===========================================================================
\section{การแยกตัวประกอบ (Factorization)}

\begin{frame}{การแยกตัวประกอบ: วิธีพื้นฐาน}
  \begin{columns}[T]
    \column{0.5\textwidth}
      \begin{explanation}[1. ดึงตัวร่วม]
        $6x^3 - 9x^2 + 3x = 3x(2x^2 - 3x + 1)$
      \end{explanation}
      \begin{explanation}[2. จับกลุ่ม (Grouping)]
        $x^3 + 2x^2 + 3x + 6 = x^2(x+2) + 3(x+2) = (x+2)(x^2+3)$
      \end{explanation}
    \column{0.5\textwidth}
      \begin{explanation}[3. ผลต่างกำลังสอง]
        $x^2 - 25 = (x-5)(x+5)$

        $4x^2 - 9y^2 = (2x-3y)(2x+3y)$
      \end{explanation}
      \begin{explanation}[4. กำลังสองสมบูรณ์]
        $x^2 + 6x + 9 = (x+3)^2$ \quad
        $x^2 - 10x + 25 = (x-5)^2$
      \end{explanation}
  \end{columns}
\end{frame}

\begin{frame}{การแยกตัวประกอบ: ผลบวก/ผลต่างกำลังสาม}
  \begin{exambox}[สูตรและตัวอย่าง]
    $a^3 + b^3 = (a+b)(a^2 - ab + b^2)$ \qquad
    $a^3 - b^3 = (a-b)(a^2 + ab + b^2)$

    $x^3-8=(x-2)(x^2+2x+4)$ \qquad
    $8x^3+27=(2x+3)(4x^2-6x+9)$
  \end{exambox}
\end{frame}

\begin{frame}{การแยกตัวประกอบ: $x^4+4$ และเทคนิคการจัดรูป}
  \begin{exambox}[เทคนิค: $x^4+4$]
    $x^4+4 = x^4 + 4x^2 + 4 - 4x^2 = (x^2+2)^2 - (2x)^2 = (x^2+2x+2)(x^2-2x+2)$
  \end{exambox}

  \begin{intuition}[แนวคิด]
    \textbf{การทำเป็นกำลังสองสมบูรณ์แล้วใช้ผลต่างกำลังสอง}:
    บวกและลบ $4x^2$ เพื่อสร้าง $(x^2+2)^2$ แล้วใช้ $A^2-B^2 = (A-B)(A+B)$
  \end{intuition}
\end{frame}

\begin{frame}{การแยกตัวประกอบโดยใช้ทฤษฎีบทตัวประกอบ}
  \begin{explanation}[แนวคิด]
    \textbf{หาราก (root)} ของพหุนามก่อน — ใช้ทฤษฎีบทตัวประกอบ — แล้วหารเพื่อลดดีกรี
  \end{explanation}

  \begin{exambox}[ตัวอย่าง: หารากและแยกตัวประกอบ]
    ทดลองแทนค่า: $P(-1) = -1-4-1+6 = 0$ $\rightarrow$ $(x+1)$ เป็นตัวประกอบ!

    หาร $P(x)$ ด้วย $(x+1)$:
    \[
    \begin{array}{c|cccc}
      -1 & 1 & -4 & 1 & 6 \\
         &    & -1 & 5 & -6 \\
      \hline
         & 1 & -5 & 6 & \boxed{0}
    \end{array}
    \]

    $P(x) = (x+1)(x^2 - 5x + 6) = (x+1)(x-2)(x-3)$
  \end{exambox}
\end{frame}

% ===========================================================================
%  SECTION 7: ทฤษฎีบททวินาม
% ===========================================================================
\section{ทฤษฎีบททวินาม (Binomial Theorem)}

\begin{frame}{สัมประสิทธิ์ทวินามและสามเหลี่ยมปัสกาล}
  \begin{explanation}[นิยาม: สัมประสิทธิ์ทวินาม]
    $\displaystyle\binom{n}{r} = \frac{n!}{r!(n-r)!}$ \quad ($0 \leq r \leq n$)
    \\ อ่านว่า ``$n$ choose $r$'' — จำนวนวิธีเลือกของ $r$ ชิ้นจาก $n$ ชิ้น
  \end{explanation}

  \begin{columns}[T]
    \column{0.45\textwidth}
      \centering
      \textbf{สามเหลี่ยมปัสกาล (Pascal's Triangle)}

      \begin{tabular}{c}
        $n=0$: \quad $1$ \\
        $n=1$: \quad $1 \quad 1$ \\
        $n=2$: \quad $1 \quad 2 \quad 1$ \\
        $n=3$: \quad $1 \quad 3 \quad 3 \quad 1$ \\
        $n=4$: \quad $1 \quad 4 \quad 6 \quad 4 \quad 1$ \\
        $n=5$: \quad $1 \quad 5 \quad 10 \quad 10 \quad 5 \quad 1$
      \end{tabular}

    \column{0.55\textwidth}
      \begin{explanation}[สมบัติ]
        \begin{itemize}
          \item $\binom{n}{0} = \binom{n}{n} = 1$
          \item $\binom{n}{r} = \binom{n}{n-r}$ \ (สมมาตร)
          \item $\binom{n}{r} = \binom{n-1}{r-1} + \binom{n-1}{r}$
          \item แต่ละแถวรวมกัน $= 2^n$
        \end{itemize}
      \end{explanation}
  \end{columns}
\end{frame}

\begin{frame}{ทฤษฎีบททวินาม}
  \begin{explanation}[ทฤษฎีบท]
    \highlight{$(a+b)^n = \displaystyle\sum_{r=0}^{n} \binom{n}{r} a^{n-r} b^r$}

    \medskip
    $= \binom{n}{0}a^n + \binom{n}{1}a^{n-1}b + \binom{n}{2}a^{n-2}b^2 + \cdots + \binom{n}{n}b^n$
  \end{explanation}

  \begin{exambox}[ตัวอย่าง: $(x+y)^3$]
    $(x+y)^3 = \binom{3}{0}x^3 + \binom{3}{1}x^2y + \binom{3}{2}xy^2 + \binom{3}{3}y^3$

    $= 1 \cdot x^3 + 3x^2y + 3xy^2 + 1 \cdot y^3$

    $= x^3 + 3x^2y + 3xy^2 + y^3$
  \end{exambox}
\end{frame}

\begin{frame}{ตัวอย่าง: ทฤษฎีบททวินาม}
  \begin{exambox}[โจทย์ 1]
    จงกระจาย $(2x - 3)^4$
  \end{exambox}

  \begin{solbox}[วิธีทำ]
    $(2x-3)^4 = \sum_{r=0}^{4} \binom{4}{r} (2x)^{4-r} (-3)^r$

    $= 16x^4 - 96x^3 + 216x^2 - 216x + 81$
  \end{solbox}
\end{frame}

\begin{frame}{ตัวอย่าง: การหาพจน์ที่ต้องการ}
  \begin{exambox}[โจทย์]
    จงหาพจน์ที่ $5$ ของการกระจาย $(x + 2)^{8}$
  \end{exambox}

  \begin{solbox}[วิธีทำ]
    พจน์ที่ $r+1$: $\binom{n}{r}a^{n-r}b^r$ \quad $\rightarrow$ \quad
    $r = 4$: \;
    $\binom{8}{4}x^{4}(2)^4 = 70 \cdot x^4 \cdot 16 = 1{,}120x^4$
  \end{solbox}

  \begin{exambox}[โจทย์ 2]
    จงหาพจน์กลางของการกระจาย $(2x - 1)^6$
  \end{exambox}

  \begin{solbox}[วิธีทำ]
    $n=6$ มี $7$ พจน์ $\rightarrow$ พจน์กลางคือพจน์ที่ $4$ ($r = 3$)

    $\binom{6}{3}(2x)^{3}(-1)^3 = 20 \cdot 8x^3 \cdot (-1) = -160x^3$
  \end{solbox}
\end{frame}

% ===========================================================================
%  SECTION 8: ทฤษฎีบทรากตรรกยะ
% ===========================================================================
\section{ทฤษฎีบทรากตรรกยะ (Rational Root Theorem)}

\begin{frame}{ทฤษฎีบทรากตรรกยะ}
  \begin{explanation}[ทฤษฎีบท]
    ถ้าพหุนาม $P(x) = a_nx^n + \cdots + a_0$ มีสัมประสิทธิ์เป็นจำนวนเต็ม
    และ $\frac{p}{q}$ เป็นรากตรรกยะในรูปเศษส่วนอย่างต่ำแล้ว:

    \highlight{$p$ หาร $a_0$ ลงตัว} \quad และ \quad \highlight{$q$ หาร $a_n$ ลงตัว}
  \end{explanation}

  \begin{exambox}[ตัวอย่าง: หารากตรรกยะ]
    $P(x) = 2x^3 - 3x^2 - 3x + 2$: \quad $a_n = 2$, $a_0 = 2$

    รากที่เป็นไปได้: $\pm 1, \pm 2, \pm \frac{1}{2}$

    ทดสอบ: $P(1) \neq 0$, \quad $P(-1) = 0 \;\checkmark$, \quad $P(2) = 0 \;\checkmark$,
    \quad $P(\frac{1}{2}) = 0 \;\checkmark$

    \highlight{$P(x) = (x+1)(x-2)(2x-1)$}
  \end{exambox}
\end{frame}

% ===========================================================================
%  SECTION 9: ตัวอย่างโจทย์
% ===========================================================================
\section{ตัวอย่างโจทย์ประยุกต์}

\begin{frame}{ตัวอย่าง: หาพหุนามจากรากที่กำหนด}
  \begin{exambox}[โจทย์]
    จงสร้างพหุนามดีกรี $3$ ที่มีสัมประสิทธิ์นำ $= 1$ และมีรากเป็น $2, -1, 4$
  \end{exambox}

  \begin{solbox}[วิธีทำ]
    ราก $= 2, -1, 4$ ดังนั้นตัวประกอบคือ $(x-2), (x+1), (x-4)$

    $P(x) = (x-2)(x+1)(x-4)$

    $= (x^2 - x - 2)(x-4)$

    $= x^3 - 4x^2 - x^2 + 4x - 2x + 8$

    $= x^3 - 5x^2 + 2x + 8$
  \end{solbox}
\end{frame}

\begin{frame}{ตัวอย่าง: หาค่า $k$ จากเงื่อนไขเศษเหลือ}
  \begin{exambox}[โจทย์]
    เมื่อหาร $P(x) = 2x^3 + kx^2 - 5x + 3$ ด้วย $(x-2)$ เหลือเศษ $7$
    จงหาค่า $k$
  \end{exambox}

  \begin{solbox}[วิธีทำ]
    จากทฤษฎีบทเศษเหลือ: $P(2) = 7$

    $2(8) + k(4) - 5(2) + 3 = 7$

    $16 + 4k - 10 + 3 = 7$

    $4k + 9 = 7$

    $4k = -2$

    \highlight{$k = -\frac{1}{2}$}
  \end{solbox}
\end{frame}

\begin{frame}{ตัวอย่าง: แยกตัวประกอบพหุนามดีกรี 4}
  \begin{exambox}[โจทย์]
    จงแยกตัวประกอบของ $P(x) = x^4 - 5x^2 + 4$
  \end{exambox}

  \begin{solbox}[วิธีทำ]
    ให้ $y = x^2$: \quad $P(x) = y^2 - 5y + 4 = (y-1)(y-4)$

    แทน $y = x^2$ กลับ: \quad $(x^2-1)(x^2-4) = (x-1)(x+1)(x-2)(x+2)$

    \highlight{$P(x) = (x-1)(x+1)(x-2)(x+2)$}
  \end{solbox}

  \begin{intuition}[เทคนิค]
    พหุนามในรูป $ax^{2n} + bx^n + c$ (quadratic in form) — ใช้การแทนค่าเพื่อลดดีกรี
  \end{intuition}
\end{frame}

% ===========================================================================
%  SUMMARY
% ===========================================================================
\section{สรุป}

\begin{frame}{สรุปเนื้อหา}
  \begin{explanation}[สิ่งที่ได้เรียนรู้]
    \begin{enumerate}
      \item \textbf{นิยาม}: $P(x) = a_nx^n + \cdots + a_0$, ดีกรี, สัมประสิทธิ์
      \item \textbf{การดำเนินการ}: บวก/ลบ (รวมพจน์เหมือน), คูณ (แจกแจง)
      \item \textbf{การหาร}: หารยาว, หารสังเคราะห์ (ตัวหาร $x-c$)
      \item \textbf{ทฤษฎีบทเศษเหลือ}: เศษจาก $(x-c) = P(c)$
      \item \textbf{ทฤษฎีบทตัวประกอบ}: $P(c)=0 \iff (x-c)$ เป็นตัวประกอบ
      \item \textbf{การแยกตัวประกอบ}: ดึงตัวร่วม, จับกลุ่ม, ผลต่างกำลังสอง,
            กำลังสองสมบูรณ์, ผลบวก/ผลต่างกำลังสาม, ใช้ทฤษฎีบทตัวประกอบ
      \item \textbf{ทฤษฎีบททวินาม}: $(a+b)^n = \sum_{r=0}^{n} \binom{n}{r}a^{n-r}b^r$,
            สามเหลี่ยมปัสกาล, $\binom{n}{r} = \frac{n!}{r!(n-r)!}$
      \item \textbf{ทฤษฎีบทรากตรรกยะ}: $\frac{p}{q}$ โดย $p \mid a_0$, $q \mid a_n$
    \end{enumerate}
  \end{explanation}
\end{frame}

\end{document}
